\section{Implementation}

This project consists of two main components:

\begin{enumerate}
    \item \textbf{Traffic-DSTG-Gen}: Dataset generation system that creates dynamic spatio-temporal graph datasets from SUMO traffic simulations
    \item \textbf{TrafficLab}: Web-based testing platform that provides real-time ETA predictions through an interactive interface
\end{enumerate}

\subsection{Development Environment}

The project was developed in a Linux environment using Docker for containerization:

\begin{itemize}
    \item \textbf{Development Container}: Ubuntu 20.04 with Python 3.10, Node.js 18+
    \item \textbf{Version Control}: Git with branching strategy for parallel development
    \item \textbf{Dependency Management}: 
    \begin{itemize}
        \item Python dependencies managed with pip (requirements.txt)
        \item JavaScript dependencies managed with npm
    \end{itemize}
    \item \textbf{Containerization}: Docker and Docker Compose for deployment
\end{itemize}

\subsection{Traffic-DSTG-Gen: Multi-Variant Dataset Generation System}

The multi-variant dataset generation system is responsible for creating synchronized data formats from a single SUMO simulation, enabling fair comparison across different model architectures for both flow forecasting and ETA prediction.

\subsubsection{Architecture}

The system follows a pipeline architecture with six main stages:

\begin{enumerate}
    \item \textbf{Simulation Stage}: Generate traffic snapshots using SUMO with configurable scenarios
    \item \textbf{Multi-Variant Extraction Stage}: Extract four synchronized variants (dynamic graph, trajectory, static graph, route segment) from the same simulation
    \item \textbf{Labeling Stage}: Generate per-snapshot labels for both flow prediction (speed, volume, occupancy) and ETA prediction (travel times)
    \item \textbf{EDA Stage}: Generate feature statistics and analysis for each variant
    \item \textbf{Conversion Stage}: Convert each variant to its native format (PyTorch Geometric for graphs, CSV/JSON for trajectories and route segments)
    \item \textbf{Validation Stage}: Verify dataset integrity, temporal synchronization, and consistent train/validation/test splits across all variants
\end{enumerate}

\subsubsection{Key Components}

\textbf{Simulation Manager} (\texttt{graph/entities.py}):
\begin{itemize}
    \item \texttt{SimManager}: Main controller for simulation state and multi-variant extraction
    \item Junction and Vehicle node abstractions
    \item Dynamic edge representations
    \item Route-aware feature encoding
    \item Multi-variant data extraction coordination
\end{itemize}

\textbf{Main Entry Point} (\texttt{main.py}):
The main entry point orchestrates the entire simulation and multi-variant extraction:
\begin{itemize}
    \item Loads SUMO network configuration
    \item Initializes traffic simulation with vehicle schedules
    \item Extracts dynamic graph snapshots at configurable intervals (default: 30 seconds)
    \item Extracts GPS trajectories for trajectory-based models
    \item Aggregates sensor readings for static graph models
    \item Extracts route segments for route-aware models
    \item Saves all variants as synchronized JSON files for further processing
\end{itemize}

\textbf{Multi-Variant Dataset Tools}:
\begin{itemize}
    \item \texttt{dataset\_creator.py}: Converts JSON snapshots to PyTorch Geometric .pt files (dynamic graph variant)
    \item \texttt{trajectory\_extractor.py}: Extracts GPS trajectory sequences from vehicle movements (trajectory variant)
    \item \texttt{static\_graph\_aggregator.py}: Aggregates sensor readings at fixed locations (static graph variant)
    \item \texttt{route\_segment\_extractor.py}: Extracts route segments with aggregated features (route segment variant)
    \item \texttt{create\_labels\_json.py}: Generates per-snapshot labels for both flow and ETA prediction
    \item \texttt{synchronization\_validator.py}: Validates temporal synchronization and consistent splits across variants
    \item \texttt{EDA.py}: Comprehensive exploratory data analysis for each variant
    \item \texttt{visualize\_graph.py}: Graph visualization utilities
\end{itemize}

\subsubsection{Multi-Variant Dataset Structure}

The generated dataset includes four synchronized variants, each tailored to different model architectures:

\textbf{Variant 1: Dynamic Graph} (for DSTRA-GNN and similar models):
\begin{itemize}
    \item \textbf{Node Features (28 dimensions)}:
    \begin{itemize}
        \item \textbf{Junction Nodes}: Zone, position, type, traffic state
        \item \textbf{Vehicle Nodes}: Speed, acceleration, position, route progress, destination, route\_left information
    \end{itemize}
    \item \textbf{Edge Features (7 dimensions)}:
    \begin{itemize}
        \item \textbf{Static Edges}: Length, lanes, speed limits
        \item \textbf{Dynamic Edges}: Current speed, demand, occupancy
    \end{itemize}
    \item \textbf{Route Awareness}: Explicit route encoding with complete remaining route per vehicle
    \item \textbf{Format}: PyTorch Geometric .pt files with temporal sequences
    \item \textbf{Labels}: Both flow (speed, volume, occupancy at junctions) and ETA (per-vehicle travel times)
\end{itemize}

\textbf{Variant 2: Trajectory} (for DeepTTE, STAD, and similar models):
\begin{itemize}
    \item \textbf{Format}: GPS trajectory sequences—ordered lists of (latitude, longitude, timestamp) tuples per trip
    \item \textbf{Features}: Trip-level metadata (origin, destination, start time, duration)
    \item \textbf{Labels}: ETA (per-trip travel times)
    \item \textbf{Format}: CSV/JSON files with trajectory sequences
\end{itemize}

\textbf{Variant 3: Static Graph} (for DCRNN, ST-GCN, and similar models):
\begin{itemize}
    \item \textbf{Format}: Aggregated sensor readings at fixed locations (junctions or road segments)
    \item \textbf{Graph Structure}: Static road network topology with fixed node and edge structure
    \item \textbf{Features}: Time-series of aggregated measurements (speed, flow, occupancy) at each sensor
    \item \textbf{Labels}: Flow prediction (speed, volume, occupancy at fixed locations)
    \item \textbf{Format}: PyTorch Geometric .pt files with static graph and time-series features
\end{itemize}

\textbf{Variant 4: Route Segment} (for DuETA and similar models):
\begin{itemize}
    \item \textbf{Format}: Trip-level records with explicit route sequences
    \item \textbf{Features}: Route represented as ordered list of road segments or waypoints, segment-level aggregated features (length, speed limit, historical travel time)
    \item \textbf{Labels}: ETA (per-trip travel times) with duration-aware categorization
    \item \textbf{Format}: CSV/JSON files with route sequences and segment features
\end{itemize}

\textbf{Synchronization and Consistency}:
\begin{itemize}
    \item All variants derived from the same underlying SUMO simulation
    \item Identical train/validation/test splits (chronological partitioning)
    \item Synchronized timestamps for temporal alignment
    \item Consistent trip identifiers across variants for cross-variant analysis
    \item Same traffic patterns and scenarios across all variants
\end{itemize}

\subsection{SUMO Integration and Multi-Variant Extraction}

The SUMO integration was implemented using the Python TraCI library to control running simulations in real-time and extract multiple synchronized variants. The integration provides:

\begin{itemize}
    \item Dynamic vehicle tracking and state extraction for all variants
    \item Traffic light state monitoring for flow prediction
    \item Network topology extraction for static graph variant
    \item Real-time route planning using Dijkstra's algorithm
    \item Graph structure generation for dynamic graph variant
    \item GPS trajectory extraction for trajectory variant
    \item Sensor aggregation at fixed locations for static graph variant
    \item Route segment extraction with aggregated features for route segment variant
    \item Temporal synchronization across all variants
    \item Consistent labeling for both flow and ETA prediction
\end{itemize}

\subsection{Model Inference Implementation}

The model inference service loads pre-trained DSTRA-GNN models and performs real-time predictions. Key implementation details include:

\begin{itemize}
    \item Model loading from checkpoint files
    \item Graph data transformation to PyTorch Geometric format
    \item Batch processing for efficient inference
    \item GPU acceleration support
    \item Result caching for performance optimization
\end{itemize}

\subsection{TrafficLab: Web-Based Testing Platform}

TrafficLab is a comprehensive web application that integrates SUMO simulation with DSTRA-GNN model inference.

\subsubsection{Backend Architecture}

\textbf{FastAPI Application} (\texttt{backend/main.py}):
\begin{itemize}
    \item RESTful API with 50+ endpoints for simulation control
    \item SUMO integration via TraCI for real-time simulation control
    \item Model inference service for ETA predictions
    \item Journey tracking and analytics
    \item CORS middleware for frontend communication
\end{itemize}

\textbf{SUMO Service} (\texttt{backend/services/sumo\_service.py}):
\begin{itemize}
    \item Manages SUMO simulation lifecycle (start, pause, stop)
    \item Real-time vehicle tracking and position updates
    \item Network topology management (junctions, edges, zones)
    \item Route planning using SUMO's internal Dijkstra algorithm
    \item Traffic state extraction for model input
\end{itemize}

\textbf{Model Inference} (\texttt{backend/models/}):
\begin{itemize}
    \item \texttt{eta\_inference.py}: Loads pre-trained DSTRA-GNN models
    \item \texttt{real\_time\_inference.py}: Real-time prediction pipeline
    \item \texttt{temporal\_dataset.py}: Handles temporal graph sequences
    \item \texttt{model\_temporal\_moe.py}: Implements Mixture-of-Experts architecture
    \item GPU acceleration support for low-latency inference
\end{itemize}

\subsubsection{Frontend Architecture}

\textbf{Vue.js Application} (\texttt{frontend/src/}):
\begin{itemize}
    \item \textbf{App.vue}: Main application shell with routing
    \item \textbf{HomePage.vue}: Landing page with research information
    \item \textbf{DemoPage.vue}: Interactive simulation demo interface
    \item \textbf{SimDemoPage.vue}: Full-featured simulation control panel
    \item \textbf{DatasetCreator.vue}: Dataset generation interface
\end{itemize}

\textbf{Key Features}:
\begin{itemize}
    \item Interactive map-based route selection
    \item Real-time traffic visualization with vehicle positions
    \item ETA prediction display with accuracy metrics
    \item Journey tracking with predicted vs actual travel times
    \item Statistical analysis dashboard with performance metrics
    \item Responsive design using Tailwind CSS
\end{itemize}

\subsubsection{Database Schema}

\textbf{PostgreSQL Database} (\texttt{backend/models/database.py}):
\begin{itemize}
    \item \texttt{Journey} table: Stores completed trips with predictions and actual times
    \item SQLAlchemy ORM for database operations
    \item Automatic schema migration support
    \item Session management with connection pooling
\end{itemize}

\subsubsection{Deployment}

\textbf{Docker Compose Setup}:
\begin{itemize}
    \item Separate containers for frontend (Vue.js) and backend (FastAPI)
    \item PostgreSQL database container with persistent storage
    \item Nginx reverse proxy for production deployment
    \item Environment-based configuration (development vs production)
    \item Health checks and automatic restart on failure
\end{itemize}

\textbf{Development Workflow}:
\begin{itemize}
    \item Hot-reload for both frontend and backend during development
    \item Shared volumes for code and data persistence
    \item Network isolation with Docker networking
    \item Log aggregation for debugging
\end{itemize}

